\chapter{Những câu hỏi làm thầy khựng lại}
\markboth{Những câu hỏi làm thầy khựng lại}{Những câu hỏi làm thầy khựng lại}

\section{Người lớn sai nhiều vậy sao cứ bắt tụi em chuẩn}
\begin{truyen}{Người lớn sai nhiều vậy sao cứ bắt tụi em chuẩn}{Classroom Talk}
\chuhoa{H}{ọc sinh: ``Người lớn cũng nóng tính, hứa rồi quên, nói một đằng làm một nẻo. Vậy sao cứ bắt tụi em phải chuẩn?''}
\textit{Student: ``Adults lose their temper too, make promises and forget them, and say one thing while doing another. So why do they keep demanding that we be perfect?''}\\[4pt]

\teacherpair{Câu đó không sai. Người lớn nhiều khi chưa làm gương đủ tốt. Nhưng cái sai của người lớn không làm cái đúng bớt đúng.}{That is not a wrong question. Adults often fail to be good examples. But the mistakes of adults do not make what is right any less right.}

\studentpair{Vậy tụi em phải nhìn vô ai?}{Then who are we supposed to look at?}

\teacherpair{Nhìn vào điều đúng, rồi cũng nhìn kỹ để sau này mình đỡ lặp lại cái dở của người lớn.}{Look at what is right, and also observe carefully so that later you can avoid repeating the worst habits of adults.}
\end{truyen}

\begin{baihoc}
Thấy người lớn sai không phải để khỏi phải tốt hơn. Là để biết mình cần sống tử tế hơn ở chỗ nào.
\textit{(Seeing adult flaws should not excuse your own. It should sharpen your sense of what to improve.)}
\end{baihoc}

\begin{ghinhoanh}
\textbf{Short English to remember:}

Adult flaws are real.

Truth does not disappear because messengers fail.
\end{ghinhoanh}

\begin{tuhocnhanh}
1. Em có đang lấy cái sai của người khác để nhẹ lỗi cho mình không? \textit{(Do you use someone else’s wrong as a discount for your own?)}

2. Điều tốt nào em vẫn muốn giữ dù từng thấy người lớn làm chưa tới? \textit{(What good value do you still want to keep even after seeing adults fail at it?)}
\end{tuhocnhanh}
\ngancach

\section{Thầy từng bất lực muốn mặc kệ chưa}
\begin{truyen}{Thầy từng bất lực muốn mặc kệ chưa}{Classroom Talk}
\chuhoa{H}{ọc sinh: ``Thầy có bao giờ dạy một lớp xong mệt quá, thấy tụi em lì quá rồi chỉ muốn mặc kệ không?''}
\textit{Student: ``Have you ever finished teaching a class, felt exhausted by how stubborn we are, and just wanted to give up on us?''}\\[4pt]

\teacherpair{Có chứ. Giáo viên cũng là người. Có lúc nói mấy lần mà không vào, lòng nản là thật.}{Of course. Teachers are human too. There are times when we explain something again and again and nothing gets through, and the discouragement is real.}

\studentpair{Vậy sao thầy còn làm?}{Then why do you keep doing it?}

\teacherpair{Vì nếu thầy mặc kệ luôn, nhiều đứa sẽ quen với cảm giác không còn ai kỳ vọng ở mình nữa.}{Because if I truly stopped caring, many students would get used to the feeling that nobody expects anything from them anymore.}
\end{truyen}

\begin{baihoc}
Sự kiên nhẫn của người dạy không vô hạn. Nhưng đôi khi nó cứu được một đứa đang sắp buông mình.
\textit{(A teacher’s patience is not unlimited, but it can sometimes rescue a student on the edge of quitting.)}
\end{baihoc}

\begin{ghinhoanh}
\textbf{Short English to remember:}

Teachers get tired too.

Staying can still change a life.
\end{ghinhoanh}

\begin{tuhocnhanh}
1. Em có từng làm ai muốn bỏ cuộc với mình chưa? \textit{(Have you ever made someone feel like giving up on you?)}

2. Nếu có người vẫn chưa bỏ em, em nên đáp lại thế nào? \textit{(If someone has not given up on you, how should you respond?)}
\end{tuhocnhanh}
\ngancach

\section{Học sinh cần bị quản hay cần được hiểu}
\begin{truyen}{Học sinh cần bị quản hay cần được hiểu}{Classroom Talk}
\chuhoa{H}{ọc sinh: ``Nhiều khi tụi em bị quản kỹ quá thấy ngộp. Nhưng nếu thả lỏng lại dễ loạn. Vậy tụi em cần bị quản hay cần được hiểu hơn?''}
\textit{Student: ``Sometimes we feel suffocated when we are controlled too closely. But if everything is loosened, things fall apart. So do students need to be managed, or understood, more?''}\\[4pt]

\teacherpair{Cả hai, nhưng theo thứ tự đúng. Hiểu để biết phải quản chỗ nào. Quản để giữ cho việc hiểu không thành nuông chiều.}{Both, but in the right order. Understanding helps us know where structure is needed. Structure keeps understanding from turning into indulgence.}

\studentpair{Nghĩa là thương không đồng nghĩa với dễ dãi?}{So caring for someone does not mean being easy on them?}

\teacherpair{Đúng. Có khi nghiêm mới là một dạng thương có trách nhiệm.}{Exactly. Sometimes being strict is a form of responsible care.}
\end{truyen}

\begin{baihoc}
Thầy cô tốt không phải lúc nào cũng mềm. Nhiều khi họ đang giữ cho mình khỏi trượt dài.
\textit{(Good teachers are not always gentle in appearance. Sometimes they are preventing a longer fall.)}
\end{baihoc}

\begin{ghinhoanh}
\textbf{Short English to remember:}

Understanding matters.

Structure matters too.
\end{ghinhoanh}

\begin{tuhocnhanh}
1. Em đang cần được cảm thông hay đang cần bị kéo lại cho đúng? \textit{(Do you need compassion, or do you need correction right now?)}

2. Em có nhầm giữa được hiểu và được chiều không? \textit{(Do you confuse being understood with being indulged?)}
\end{tuhocnhanh}
\ngancach

\section{Có đứa im không phải ngoan, mà là hết muốn nói}
\begin{truyen}{Có đứa im không phải ngoan, mà là hết muốn nói}{Classroom Talk}
\chuhoa{H}{ọc sinh: ``Mấy bạn ít nói, ít phá có chắc là ổn không thầy? Hay có đứa im vì chẳng còn muốn mở miệng nữa?''}
\textit{Student: ``Are the quiet students who never cause trouble always actually okay? Or are some of them silent because they no longer want to speak at all?''}\\[4pt]

\teacherpair{Em hỏi đúng chỗ đau của nhiều lớp. Không phải học sinh im nào cũng ngoan. Có em im vì sợ, vì mệt, vì thấy nói cũng chẳng ai nghe.}{You are touching a painful truth in many classrooms. Not every quiet student is simply well-behaved. Some stay silent because they are afraid, exhausted, or convinced that nobody will listen anyway.}

\studentpair{Vậy thầy cô làm sao biết?}{Then how can teachers tell the difference?}

\teacherpair{Phải nhìn lâu hơn điểm số và hành vi bề mặt. Nhiều đứa ngoan quá trên giấy nhưng đang mòn dần ở bên trong.}{We have to look beyond grades and surface behavior for a longer time. Many students look perfectly good on paper while slowly wearing down inside.}
\end{truyen}

\begin{baihoc}
Đừng chỉ chú ý những đứa ồn. Nhiều khi đứa cần cứu nhất lại là đứa im nhất.
\textit{(Do not only notice the loud students. The quietest one may need help the most.)}
\end{baihoc}

\begin{ghinhoanh}
\textbf{Short English to remember:}

Silence is not always peace.

Quiet students need attention too.
\end{ghinhoanh}

\begin{tuhocnhanh}
1. Em có đang im vì bình an hay vì cạn năng lượng? \textit{(Are you quiet because you are calm, or because you are drained?)}

2. Em có người lớn nào em có thể nói thật không? \textit{(Do you have an adult you can speak honestly with?)}
\end{tuhocnhanh}
\ngancach